\section{Discussion}
\label{sec:discussion}

\subsection{Why do reminders help only with a contract?}
Reminders on their own do not change behavior; they work only as a complement to the contract. A reminder tells a participant that today is a gym day, but without a contract skipping the visit costs nothing. With a contract, the reminder arrives while the stake is at risk, which may be why the combination works better than either tool alone. The pattern is also consistent with limited attention to a distant target: a contract holder who forgets the target until the last week may no longer be able to reach it.

\subsection{What does take-up reveal?}
Take-up need not reveal sophistication about present bias. \citet{carrera2018plans} find that between 27 and 34 percent of gym members accept a contract that pays them only if they visit the gym less often, which no member trying to exercise more should want. In our setting, a third of those offered a contract signed it, and 41 percent of the signers missed the target and lost their stake. Some of these signers may have overestimated their future self-control, as the partial naivet\'e documented by \citet{acland2015} would predict.

\subsection{Costs}
Administering the contracts, including staff time and card fees, cost 9.80 dollars per employee offered a contract. Since the offer raised visits by 0.31 per week over eight weeks, or 2.48 visits per employee offered a contract, each additional visit cost about 2.95 dollars. The forfeited stakes went to the food bank and are a transfer rather than a cost.

\subsection{Scope and limitations}
Our participants had already bought a membership, so the estimates speak to people who have taken a first step toward exercising more. Our data also have limits. We observe visits but not what people do at the gym, the sample comes from a single employer, and the follow-up ends twelve weeks after the contracts expired.
